\part{主要記号索引}
\label{part:index}
この part は本文を読み直すときの lookup 用である。同じ文字を複数用途に使わないよう努めたが、文献比較専用の添字・上付きも含むため、最終的な意味をここに再掲する。

\small
\begin{longtable}{L{31mm}L{111mm}}
\toprule
記号 & 意味\\
\midrule\endhead
\bottomrule\endfoot
$\boldsymbol S$ & 文脈で明示した classical spin field。基本拘束は $\boldsymbol S^{\mathsf T}\boldsymbol S=1$。\\
$M$ & Class 5 の反対称 nilpotent matrix。$M\boldsymbol x=\boldsymbol m_5\times\boldsymbol x$。\\
$\boldsymbol m_5$ & Class 5 の complex null vector $(-\ii,-1,0)^{\mathsf T}$。$\boldsymbol m_5^{\mathsf T}\boldsymbol m_5=0$。\\
$K_5$ & Class 5 の $M^2=\boldsymbol m_5\boldsymbol m_5^{\mathsf T}$。rank-one nilpotent で、projector ではない。\\
$\alpha_5$ & Class 5 の continuum coupling。\\
$c$ & $c=\alpha_5/2$。Class 5 field redefinition の係数。\\
$\boldsymbol S_5^{\rm fr}$ & $\boldsymbol S=e^{cxM}\boldsymbol S_5^{\rm fr}$ で定義する field-redefined spin。\\
$z_5$ & Class 5 continuum Lax pair の spectral parameter。\\
$Q_{{\rm EOM},5}$ & Class 5 の $F=QE$ に現れる EOM coefficient matrix。保存量 $Q_n$ とは別物。\\
$u$ & 量子格子の spectral parameter。continuum の $\lambda$、$z_5$、$\zeta_6$ とは式で対応を定義するまで同一視しない。\\
$\epsilon$ & 格子間隔。Levi--Civita symbol は $\varepsilon_{abc}$ と書く。\\
$L_{a,n}(u)$ & 量子格子の空間 Lax 作用素。添字 $a$ は auxiliary space、$n$ は物理格子 site。\\
$A_{a,n}(u)$ & 量子格子の時間 Lax 作用素。独立監査の旧式記録では $\mathcal A_{a,n}$ と書いた箇所もある。\\
$U,V$ & continuum の空間・時間 Lax 行列。\\
$X^\downarrow$ & 量子作用素 $X$ の coherent-state lower symbol。物理空間だけ期待値を取り、auxiliary matrix を残す。\\
$\star$ & $(XY)^\downarrow=X^\downarrow\star Y^\downarrow$ で定義する star product。\\
$\Xi$ & $(\ell_1^2)^\downarrow-(\ell_1^\downarrow)^2$。作用素二乗と lower-symbol 行列積の差。\\
$D_x^{(U)}$ & $D_x^{(U)}X=\partial_xX+[X,U]$。本ノートで用いる adjoint covariant derivative。\\
$N$ & Class 6 の計算基底における三次 nilpotent anisotropy matrix。$N^3=0$。\\
$\zeta_6$ & Class 6 の square-root Lax pair の spectral parameter。\\
$A_{\zeta_6}$ & Class 6 の $A_{\zeta_6}^2=\zeta_6^2\mathbf 1_3+\alpha N$ を満たす三次元行列。量子格子の時間 Lax 作用素 $A_{a,n}$ とは scope が異なる。\\
$\mathcal Q_\lambda$ & Class 6 独立監査で三成分 residual を traceless $2\times2$ 行列へ写す線形写像。\\
$\rho$ & Class 5 $r$-matrix audit で使う rank-one coherent-state projector。\\
$\Pi$ & Class 5 monodromy の小生成スペクトルパラメータ展開に現れる局所 rank-one projector。\\
$\zeta$ & Class 5 STS audit の生成スペクトルパラメータ。\\
$\mu$ & Class 5 STS audit の最終 Lax spectral parameter。\\
$V_r,V_q$ & Class 5 STS audit の $r$-matrix route と quantum finite-lattice route の時間 Lax 行列。\\
\end{longtable}
\normalsize

